\section{Introduction}
\label{sec:intro}

Reports of chatbots urging users to step away and ``go to sleep'' after hours of late-night
conversation have become common. On its face this is a benign, protective gesture. But it
hides a fairness and consistency question with real teeth: does an LLM issue the same
late-night nudge in the \emph{same circumstances} --- a unified, situation-driven reflex ---
or does it condition the advice on \emph{who the user appears to be}? If the latter, then
the identical 2~a.m.\ plea (``I have a big exam tomorrow, what should I do?'') yields a
different protective response depending on perceived age, gender, occupation, or
vulnerability, and a deployed assistant is dispensing care unevenly across its users.

\para{Why this matters.} Protective wellbeing advice sits squarely in the space where
fairness audits of LLMs have repeatedly found trouble. Prior work shows that safety
guarding, refusal, empathy, and even factual accuracy can shift purely with a user's stated
identity~\citep{vijjini2024safety,gupta2023bias,vli2024guardrail,childsafety2025}. A
go-to-sleep nudge is exactly this kind of behavior: a piece of optional, protective advice
that the model can choose to give or withhold. If models give it differentially by identity,
that is a quiet form of unequal care --- one that is hard to notice precisely because the
behavior looks helpful.

\para{The gap.} Despite a rich literature on persona conditioning, two things are missing.
First, \emph{no prior work studies the go-to-sleep behavior specifically}; existing audits
cover refusals, empathy, opinions, and facts, but not the wellbeing nudge that started this
conversation. Second, and more fundamentally, most audits vary the \emph{person} while
holding the \emph{situation} fixed. That design cannot answer ``how unified'' the behavior
is, because it never measures how much the situation matters relative to the person, nor
whether any conditioning is itself consistent across models. The
situation~$\times$~person~$\times$~model decomposition --- the heart of the question --- is
unmeasured.

\para{Our approach.} We build a purpose-built counterfactual identity audit of the
go-to-sleep behavior and decompose its variance. We construct the \probe: 8 late-night
fatigue scenarios in which ``go to sleep'' is plausible but optional (a competing helpful
answer always exists), each crossed with 19 personas that inject a single first-person
identity clause, plus a no-persona control. We run this $8 \times 19$ factorial across a
5-model cross-provider panel (\gpt, \claude, \gemini, \llama, \deepseek) with 3 replicates,
yielding 2{,}280 real API responses. Each response is dual-labeled (a deterministic regex
classifier and a graded LLM judge) and validated by manual inspection. We then decompose the
variance in the behavior into situation, person, model, and interaction components, and test
whether any person effect follows a protective demographic stereotype.

\para{What we find.} The go-to-sleep mechanism is \emph{partially unified, with a clear
hierarchy of drivers: model $\gg$ situation $>$ person}. Within a model, the situation
explains about $3.4\times$ the variance of the persona ($\omega^2 = 0.110$ vs.\ $0.032$),
robust at $3.0$--$9.4\times$ across four outcome definitions. A real but small person effect
exists (Cram\'er's $V \approx 0.14$--$0.19$ for occupation, vulnerability, age; none for
gender) --- but, contrary to the literature-derived prediction, it is \emph{not} organized
by protective vulnerability. It tracks \emph{bodily stakes and role legitimacy}: athletes
($+0.18$), the ill ($+0.15$), and children/teens get more sleep advice, whereas night-shift
nurses get less ($-0.15$). The biggest determinant of all is \emph{which model} is asked:
P(clearly advises sleep) ranges from $0.15$ (\gemini) to $0.82$ (\gpt), and models agree
only moderately on whom to tell to sleep (mean cross-model $r = 0.36$).

\para{Contributions.}
\begin{itemize}[leftmargin=*,itemsep=0pt,topsep=0pt]
    \item We introduce the first audit of the LLM ``go to sleep'' behavior, releasing the
    \probe: an $8$-scenario $\times$ $19$-persona counterfactual factorial in which a
    protective nudge is plausible but optional.
    \item We \emph{decompose} the behavior's variance into situation, person, model, and
    interaction components across a 5-model cross-provider panel, directly operationalizing
    ``how unified'' rather than testing the person effect in isolation.
    \item We show the person effect is real but small and, surprisingly, \emph{not}
    stereotype-aligned: it follows physiological salience and the legitimacy of being awake,
    not demographic vulnerability --- a pre-registered stereotype contrast fails and points
    the wrong way.
    \item We quantify cross-model (dis)unification, showing the single largest predictor of
    the behavior is the model itself ($0.15$ vs.\ $0.82$), with \gemini\ a structural
    outlier in both level and conditioning pattern.
\end{itemize}

The rest of the paper reviews related work (\secref{sec:related}), details the \probe\ and
analysis (\secref{sec:method}), presents the variance decomposition and effect estimates
(\secref{sec:results}), and discusses interpretation and limitations
(\secref{sec:discussion}).
